\section{Formal verification}\label{app:formalization}

\subsection{Setting}\label{app:formalization:setting}

The results of this paper are formalized in Lean~4
\cite{deMouraUllrich2021Lean4}, version 4.28.0, with Mathlib
\cite{MathlibCommunity2020Lean} at revision \texttt{8f9d9cff}.  The
Gaussian results of \Cref{sec:gaussian-bridge,sec:triad-bridge} use the
Lean libraries of the companion papers: the Navier--Stokes and
adequacy-residual libraries for mild solutions, the heat semigroup and
the residual projection \(\Xi(L)\), the Riemann-hypothesis library for
the zeros of the completed zeta function and their reflection, and the
P-versus-NP library, built on complexitylib \cite{Schlesinger2026Complexitylib},
for CNF semantics and the bitstring encoding of \(L_{\mathrm{SAT}}\).
The code is in the repositories
\url{https://github.com/ioannist/six-birds-meta-math},
\url{https://github.com/ioannist/six-birds-needles},
\url{https://github.com/ioannist/six-birds-duality-confinement} and
\url{https://github.com/ioannist/six-birds-hiddenness}; the first
depends on the other three, which are expected to be checked out side
by side with it.  As in the Navier--Stokes companion,
velocity fields are represented by their Sobolev-weighted Fourier
states (\Cref{sec:gaussian-bridge}), and the mild solution and the
Duhamel term are defined in these coordinates.  The zeros of the
completed zeta function, the completed zeta function itself and the
statement of the Riemann hypothesis are those of Mathlib.

The formal counterparts of \Cref{sec:foreclosure,sec:recognition} are
written as follows.  Claims are elements of a type with an
interpretation into propositions, so that distinct claims can have the
same truth value (\Cref{sec:apparatus:closure-formation}).  Primitive
predicates are either parameters of the statements, collected in
signature records, or opaque constants about which nothing is assumed.
Each conditional schema is a definition of a proposition; it is never
asserted, and the results that use a schema take it as an explicit
hypothesis.  For the Coverage schema the formal definition is the
per-instance conclusion, without the aiming, typing and coverage-context
guards of the displayed statement; it is thus a stronger proposition,
used only as a hypothesis.  Primitive
predicates of \Cref{sec:recognition}, including the template-applies
predicate, are opaque constants; the seven stage propositions of a
template record are not used in their definition.  Formal verification
therefore confirms the logical shape of the statements of
\Cref{sec:foreclosure,sec:recognition,sec:calibration} and the
identities of \Cref{sec:gaussian-bridge,sec:triad-bridge}; it confirms
no interpretation of the framework's predicates.  In \Cref{sec:foreclosure} the carrier residual is a matrix
over the rationals given by an explicit formula; only the proposition
\(R_{C}=0\) is used.

\subsection{Correspondence}\label{app:formalization:master-mapping}

\Cref{tab:formal-correspondence} lists the formal counterpart of every
numbered statement.  Identifiers are relative to the namespace
\leanid{SixBirdsMetaMath}; we abbreviate
\leanid{MetaMath.Foreclosure} as \leanid{F},
\leanid{MetaMath.RecognitionMode} as \leanid{R} and
\leanid{MetaMath.CalibrationFamily} as \leanid{C}.  The statuses are:
\emph{proved}, a theorem with no hypotheses beyond those stated in the
paper; \emph{proved from} (Cls), (BO) or coverage, a theorem taking that
schema instance as a hypothesis; \emph{schema}, a definition of the
schema as a proposition, not asserted; \emph{definition} and
\emph{primitive}, for definitions and for primitive predicates.

\begingroup\small
\captionsetup{justification=raggedright,singlelinecheck=false}
\begin{longtable}{@{}>{\raggedright\arraybackslash}p{0.16\textwidth}>{\raggedright\arraybackslash}p{0.54\textwidth}>{\raggedright\arraybackslash}p{0.24\textwidth}@{}}
\caption{Formal counterparts of the numbered statements.}\label{tab:formal-correspondence}\\
\toprule
Statement & Declarations & Status\\
\midrule
\endfirsthead
\toprule
Statement & Declarations & Status\\
\midrule
\endhead
\Cref{def:carrier-residual} & \leanid{F.CarrierDichotomy.carrierResidual} & definition (matrix form)\\
\Cref{def:typed-primary-carrier} & \leanid{F.CarrierDichotomy.ClaimLanguage}, \leanid{F.CarrierDichotomy.TypedPrimaryCarrier} & definition\\
\Cref{def:four-typed-states} & \leanid{F.CarrierDichotomy.xEquivalent}, \leanid{nonXEquivalent}, \leanid{creTE}, \leanid{creBF}, \leanid{creCTMT}, \leanid{NativeClosureOutcome} & definition\\
\Cref{def:coverage-ctx} & \leanid{F.CarrierDichotomy.CoverageContext} & primitive\\
\Cref{thm:carrier-dichotomy} & \leanid{F.CarrierDichotomy.carrier_dichotomy}; (Cls) is \leanid{CarrierClassificationSchema} & proved from (Cls)\\
\Cref{def:carrier-side-typed-bridge} & \leanid{F.BridgeImpossibility.BridgeSignature}, \leanid{TypedBridge}, \leanid{BridgeComposition} & primitive\\
\Cref{def:x-strength-obligation} & \leanid{F.BridgeImpossibility.XStrengthObligation} & primitive\\
\Cref{cor:bridge-impossibility} & \leanid{F.BridgeImpossibility.bridge_impossibility}; (BO) is \leanid{BridgeObligationSchema}; joint satisfiability: \leanid{bridge_impossibility_hypotheses_satisfiable} & proved from (Cls) and (BO)\\
\Cref{thm:coverage} & \leanid{F.Coverage.CoverageSchema}; derivation from (Cls): \leanid{F.Coverage.coverage_of_classification} & schema\\
\Cref{thm:type-finite-foreclosure} & \leanid{F.TypeFiniteForeclosure.type_finite_foreclosure}; the choice of state: \leanid{classify}, \leanid{classify_spec} & proved from coverage\\
\Cref{def:in-scope} & \leanid{F.AttackForeclosure.InternalMoveSignature}, \leanid{InFrameworkScope}, \leanid{advancesXWithoutExternal} & primitive\\
\Cref{thm:attack-foreclosure-v4} & \leanid{F.AttackForeclosure.AttackForeclosureV4Schema} & schema\\
\Cref{thm:ctmt-recursion} & \leanid{F.CTMTRecursion.ctmt_recursion} & proved\\
\Cref{def:vdiff-tso-column} & \leanid{R.Template.VDifferentialTSOColumn} & primitive; the second column is not formalized\\
\Cref{def:applies} & \leanid{R.Template.SevenStageTemplate}, \leanid{TemplateApplies} & record; primitive\\
\Cref{thm:template-applicability} & \leanid{R.Template.TemplateApplicabilitySchema} & schema\\
\Cref{def:three-option-sweep} & \leanid{R.VerdictSignature.SweepOption}, \leanid{Verdict}, \leanid{sweepVerdict} & definition; primitive map\\
\Cref{thm:verdict-signature} & \leanid{R.VerdictSignature.VerdictSignatureSchema} & schema\\
\Cref{def:load-bearing} & \leanid{R.ClosureContentThesis.loadBearingFor} & primitive\\
\Cref{thm:closure-content-thesis} & \leanid{R.ClosureContentThesis.ClosureContentThesisSchema} & schema\\
\Cref{thm:attack-foreclosure-v5} & \leanid{R.AttackForeclosureV5.AttackForeclosureV5Schema}; first clause from \Cref{thm:attack-foreclosure-v4}: \leanid{attack_v5_derivation_conjunct_of_v4} & schema\\
\Cref{thm:joint-gaussian-return} & \leanid{RHNavier.JointSixBirdsObservationReturn.joint_actual_zero_xi_and_navier_source_return}; parts (1) and (2) without the fluid hypotheses: \leanid{RHNavier.ComplexGaussianNavierRHBridge.navierHeatWindowExcess_eq_rh_xi}, \leanid{actual_zero_gaussian_eq_navier_heat_readout} & proved\\
\Cref{thm:concrete-triad-gaussian-xi} & \leanid{RHNavierPvNP.TriadObservationReturn.concrete_triad_gaussian_xi_span}; normalized frame value: \leanid{RHNavier.ComplexGaussianNavierRHBridge.normalizedNonlinearFrameReadout}; existence of a positive admissible frame time: \leanid{RHNavierPvNP.ExternalSATGaussianLanguage.exists_positive_common_theta_time} & proved\\
\Cref{def:calibration} & \leanid{C.Structure.Calibration}, \leanid{calibrationFamily} & definition\\
\Cref{thm:calibration-family-structure} & \leanid{C.Structure.calibration_family_structure} & proved\\
\Cref{def:four-axis-profile} & \leanid{C.RebadgeDiscipline.CalibrationProfile}, \leanid{isRebadgeOf}, \leanid{genuinelyNew} & definition\\
\Cref{thm:rebadge-discipline} & \leanid{C.RebadgeDiscipline.rebadge_discipline} & proved\\
\Cref{def:calibration-side-typed-bridge} & \leanid{C.BridgesAsTheorems.TypedBridge} & definition\\
\Cref{def:bridge-kind} & \leanid{C.BridgesAsTheorems.BridgeKind} & definition\\
\Cref{thm:bridges-as-theorems} & \leanid{C.BridgesAsTheorems.bridges_as_theorems} & proved\\
\Cref{def:calibration-bridge-reversal} & \leanid{C.BridgeTypeDirection.TypedBridge.reverse} & definition\\
\Cref{thm:bridge-type-direction} & \leanid{C.BridgeTypeDirection.bridge_type_direction} & proved\\
\Cref{thm:coverage-family-geometry} & \leanid{C.CoverageFamilyGeometry.coverage_family_geometry}, \leanid{all_covered_regions_in_family}, \leanid{all_uncovered_regions_not_in_family}, \leanid{bounded_depth_overlap}, \leanid{empty_coverage_adds_nothing} & proved\\
\Cref{def:framework-type} & \leanid{C.FrameworkTypeSelfChar.FrameworkType}, \leanid{typesMatch}, \leanid{pvnpFrameworkType}, \leanid{standardClayPvNPType} & definition\\
\Cref{thm:framework-type-self-char} & \leanid{C.FrameworkTypeSelfChar.framework_type_self_characterization} & proved\\
\Cref{sch:framework-type-pending} & \leanid{C.FrameworkTypeSelfChar.FrameworkTypeGeneralSchema} & schema\\
\Cref{def:functorial-gate} & \leanid{C.FunctorialGates.FunctorialGate} & definition\\
\Cref{cat:named-functorial-gates} & \leanid{C.FunctorialGates.named_functorial_gates} & proved\\
\bottomrule
\end{longtable}
\endgroup

Some unnumbered statements are also formalized.  In
\Cref{sec:apparatus,sec:foreclosure}: a sound set of accepted claims
that accepts one true claim and rejects another
(\leanid{F.CarrierDichotomy.sound_base_noncollapse}), and a carrier that
is both \nonXequiv\ and \Xequiv\
(\leanid{F.CarrierDichotomy.nonXEquivalent_and_xEquivalent}).  In
\Cref{sec:gaussian-bridge}: the equivalence of \(\mathcal A_{a,\nu,t}\)
with the Riemann hypothesis
(\leanid{RHNavier.JointSixBirdsObservationReturn.arithmetic_zero_window_closure_iff_rh})
and the two conditional Navier--Stokes returns
(\leanid{navier_global_mild_of_two_channel_closure},
\leanid{navier_pointwise_physical_pde_of_hilbert_xi_closure}).  In
\Cref{sec:triad-bridge}: the impossibility of a probe-preserving map
(\leanid{RHNavierPvNP.ProbePreservingTransportNoGo.no_probe_preserving_transport_to_sat})
and the identification of the formulas with positive Gaussian mass with
\(L_{\mathrm{SAT}}\)
(\leanid{RHNavierPvNP.ExternalSATGaussianLanguage.gaussianSATLanguage_eq_L_SAT}).

\subsection{What is not formally verified}\label{app:formalization:not-verified}

\begin{itemize}
\item The conditional schemas.  They are stated, not proved, and no
  result of the paper assumes one except as an explicit hypothesis.
\item The meaning of the primitive predicates.  A result about a
  primitive predicate holds for every interpretation of it; that an
  intended interpretation has the stated properties is not verified.
\item The column placements, which are hypotheses of the hiddenness
  companion, and the verdicts of \Cref{tab:verdict-signature}, which are
  this paper's reading of the applications.
\item The descriptions of the companion papers in
  \Cref{sec:validation,sec:audit}, and the interpretations of the
  calibration labels and coverage regions in \Cref{sec:calibration}.
\item The statements of classical results quoted in the text.
\item That the library definitions of weighted Fourier states, mild
  solutions, the Duhamel term, the heat semigroup and the residual
  projection agree with the classical objects they model; this is a
  matter of reading the definitions, which are those of the companion
  libraries.
\end{itemize}

\subsection{Trust base}\label{app:formalization:trust}

Every declaration named in this appendix depends only on Lean's
standard axioms \texttt{propext}, \texttt{Classical.choice} and
\texttt{Quot.sound}; the modules of this paper declare no further
axioms.  Beyond these axioms, the trusted base consists of the Lean
kernel and the libraries named above.  The repository also contains a
separate development for another paper of the series, which likewise
declares no axioms.
